\begin{longtable}{@{}l p{12.5cm}@{}}
  \caption{Skills per-instance answer comparison ($-25\% < \Delta < 0$) for the 15 most divergent items: same question, each model's answer and reasoning.}\label{tab:cmp_skills_neg0} \\
  \toprule
  \endfirsthead
  \toprule
  \endhead
  \multicolumn{2}{@{}l}{\textbf{15228} \quad $\Delta = -20\%$} \\
  \textbf{Question} & Based on an analysis of a user's educational values and the compatibility with existing higher education biology programs, describe the key components you would include in an environmental analysis report for an educational facility. In your answer, list and define each component, explain how it connects to the user's educational values and the existing education system, and summarize how these components together inform recommendations for facility characteristics. \\
  \textbf{Reference} & An environmental analysis report for an educational facility---grounded in the user's educational values and compatibility with existing higher education biology programs---should include the following key components: \newline 1. Executive Summary: A concise overview that summarizes the user's core educational values, the degree of compatibility with existing biology programs, major findings about the facility environment, and top recommendations. This connects to the user's values by foregrounding priorities (e.g., inquiry-based learning, lab safety, interdisciplinary research) so stakeholders immediately see alignment. \newline 2. Context and Background: Definition of the institutional context (type of higher education institution, program accreditation requirements, local/regional educational policies) and a description of the user's stated educational values (e.g., student-centered pedagogy, experiential learning, sustainability). This section establishes the baseline for compatibility with existing systems by mapping requirements and values. \newline 3. Stakeholder Analysis: Identification and brief description of stakeholders (students, faculty, administrators, technical staff, community partners) and their priorities. Define how each stakeholder's expectations reflect or conflict with the user's educational values---for example, faculty emphasis on research may align with values promoting inquiry, while limited technical staff may constrain hands-on lab experiences. \newline 4. Physical Infrastructure Assessment: Definition and description of facility characteristics---laboratories (wet/dry labs), classrooms, greenhouses, fieldwork spaces, specimen storage, safety equipment, and accessibility. For each element, note current condition, capacity, and suitability for biology curricula. Link this to the user's values by assessing whether spaces support active learning, collaborative research, and required biosafety levels. \newline 5. Educational Resources and Curriculum Compatibility: Inventory and definition of curricular needs (core biology courses, lab modules, elective specializations) and resources (equipment, teaching materials, digital platforms). Explain how these map to the existing education system's standards and the user's pedagogical priorities---for instance, whether existing curricula allow for inquiry-based lab modules and how facility resources support them. \newline \textit{[\ldots\ truncated \ldots]} \newline 11. Implementation Considerations and Next Steps: A brief summary of practical steps---prioritization of upgrades, estimated sequencing (short-, medium-, long-term), and stakeholders responsible for actions. This ensures the report connects analysis to actionable planning while respecting existing institutional processes. \newline Together, these components provide a comprehensive picture: definitions and inventories clarify what exists; stakeholder and curriculum analyses show how well the environment supports the user's educational values; infrastructure, safety, and staffing assessments reveal constraints; and the SWOT plus recommendations translate analysis into targeted facility characteristics. The synthesis enables decision-makers to ensure the facility design and upgrades will both reflect the user's pedagogical priorities (e.g., experiential, research-led biology education) and remain compatible with the existing higher education system's curricular and regulatory frameworks. \\
  \midrule
  \textbf{Kimi-K2.6} & score 8/10 \\
  & [7-8] The report reflects the context well, addressing key issues with a logical structure and clear connections to educational values and the existing system. It provides specific definitions and a strong summary, but lacks some detailed evidence, explicit data sources, and a concrete step-by-step action plan, keeping it just shy of the top tier. \\
  \textbf{Innospark-1T} & score 6/10 \\
  & Let me think through this carefully. \newline The question/task: Based on an analysis of a user's educational values and the compatibility with existing higher education biology programs, describe the key components you would include in an environmental analysis report for an educational facility. In your answer, list and define each component, explain how it connects to the user's educational values and th \newline My reasoning: Environmental Analysis Report for Educational Facility \newline Thinking Process \newline To create a comprehensive environmental analysis report for an educational facility, I need to systematically examine both external and internal factors that would influence the facility's design and function. This analysis must align with educational values and integrate with existing higher education biology programs. Key Components of an Environmental Analysis Report \newline 1. Demographic Analysis \newline \textit{[\ldots\ truncated \ldots]} \newline | Economic Factors | Phased implementation, cost-effective solutions, funding strategies | \newline By synthesizing insights from all components, the environmental analysis report can provide evidence-based recommendations for facility characteristics that align with educational values while ensuring compatibility with existing higher education biology programs. \\
  \specialrule{1.2pt}{8pt}{8pt}
  \multicolumn{2}{@{}l}{\textbf{52688} \quad $\Delta = -20\%$} \\
  \textbf{Question} & You are a teaching assistant supporting an instructor who must provide personalized learning paths after a group legal skills practicum. Based on aggregated data showing that some students prioritize professional ethics while others prioritize case outcome efficiency, and additional patterns such as differences in question frequency, common misconceptions about evidentiary rules, and varying confidence levels, design an individualized learning path for a student who (a) frequently asks clarifying questions, (b) demonstrates correct procedural steps but shows inconsistent application of ethical considerations, and (c) reports low confidence in courtroom interactions. In your answer, (1) prioritize the affective values that should guide the learning path and resolve any conflicts between values (for example, ethical adherence vs. outcome-driven tactics); (2) outline 4 concrete learning activities (including one collaborative, one reflective, one practice-based, and one feedback-loop activity) ordered by priority; and (3) explain how each activity addresses the student's affective development (attitudes, values, confidence) and how you would adjust the path if the student begins to favor outcome efficiency over ethics after the first two weeks. \\
  \textbf{Reference} & Prioritization of values and conflict resolution \newline Primary values to prioritize are professional integrity (ethical adherence), client-centered responsibility (serving client interests within ethical bounds), and confidence-building through competence. These are ordered with professional integrity first because legal practice depends on adherence to ethical rules as a nonnegotiable foundation; client-centered responsibility is second because effective advocacy serves the client's interests but must operate within ethical constraints; confidence and efficiency are third but important for practical performance. To resolve conflict between ethics and outcome-driven tactics, the guiding principle is: choose ethically permissible strategies that maximize client outcomes. Where outcome efficiency would push toward ethically dubious tactics, require the student to adopt alternative efficient strategies that are ethically sound. Frame tensions as professional dilemmas to be managed, not as competing absolutes. \newline Learning activities (ordered by priority) \newline 1) Guided reflective case journal (Reflective) --- Priority 1 \newline - Activity: Student maintains a structured weekly journal of three recent decisions from practicum: (i) the procedural choice made, (ii) any ethical questions encountered, and (iii) alternative ethically permissible tactics that could improve outcomes. Each entry must include the student's emotional response and confidence rating for the interaction. \newline - Affective impact: Encourages internalization of professional integrity by prompting explicit reflection on ethics, strengthens value alignment by making trade-offs explicit, and builds metacognitive awareness that reduces anxiety. Documenting emotions and confidence provides baseline measures to track growth. \newline \textit{[\ldots\ truncated \ldots]} \newline Overall rationale and expected outcomes \newline This path emphasizes internalizing ethical values while maintaining a pragmatic focus on client outcomes. Reflection secures value alignment; practice builds competence and confidence; collaborative analysis develops normative judgment; and rapid feedback ensures behavior change. The ordering prioritizes internalization (reflective work) first because affective growth requires reflection on values, then moves to applying those values under pressure. If the student shifts toward outcome-only thinking, the targeted adjustments re-center ethics through reflection, example, and accountability rather than punishment, thereby fostering a stable professional identity that supports both ethical adherence and effective advocacy. \\
  \midrule
  \textbf{Kimi-K2.6} & score 9/10 \\
  & [9-10] The response provides highly sophisticated individual path suggestions perfectly tailored to the student's profile. Values are clearly prioritized with a strong conflict resolution framework. The four activities are concrete, correctly categorized, appropriately ordered, and explicitly linked to the student's affective needs. The adjustment plan is detailed, feasible, and directly addresses the specific risk of shifting priorities, demonstrating outstanding personalization suitability and expected learning effect. \\
  \textbf{Innospark-1T} & score 7/10 \\
  & Let me think through this carefully. \newline The question/task: You are a teaching assistant supporting an instructor who must provide personalized learning paths after a group legal skills practicum. Based on aggregated data showing that some students prioritize professional ethics while others prioritize case outcome efficiency, and additional patterns such as differences in question frequency, common mi \newline My reasoning: \textless{}analysis\textgreater{} \newline Let me analyze this problem systematically: \newline 1. Student Profile Analysis: \newline - Frequently asks clarifying questions indicates metacognitive awareness and willingness to seek help \newline \textit{[\ldots\ truncated \ldots]} \newline This individualized path addresses the student's specific needs: the clarifying questions indicate readiness for reflective practice; the procedural competence provides foundation for ethical integration; and the low confidence is addressed through scaffolded, supported experiences that build both competence and ethical grounding. \newline \textless{}/answer\textgreater{} \\
  \specialrule{1.2pt}{8pt}{8pt}
  \multicolumn{2}{@{}l}{\textbf{49079} \quad $\Delta = -20\%$} \\
  \textbf{Question} & You are a teaching assistant supporting an instructor who will receive a personalized learner summary and suggested learning path based on analysis of an individual student's response patterns, question frequency, and demonstrated understanding in a middle-school math unit on proportional relationships. Using the affective (valuing) domain at the Valuing level, analyze how you would prioritize which instructional strategies and motivational supports the instructor should offer this student. In your response, (1) identify three specific learner signals (from response patterns, question frequency, or understanding) that would influence your priorities, (2) explain the value you assign to each signal in shaping the learning path, and (3) describe concrete instructional strategies and supports you would recommend for each signal to foster the student's engagement and sense of ownership over learning. \\
  \textbf{Reference} & Learner signals I would prioritize: (a) frequent low-confidence indicators on correct answers (student answers correctly but marks low confidence or submits many clarification questions), (b) concentrated incorrect responses on scaling word problems while other proportional tasks are mostly correct, and (c) a pattern of long response times and few attempts on open-ended reflective prompts (suggesting low willingness to self-express). I value these signals as follows and recommend corresponding strategies: \newline 1) Frequent low-confidence despite correct answers \newline Value: This signal indicates the student values accuracy but lacks self-efficacy; their knowledge is emerging but not internalized. Increasing confidence will help them take productive risks and apply skills in new contexts. \newline Recommended supports: Provide brief, scaffolded justification prompts after correct responses (e.g., "Explain in one sentence why your answer makes sense"). Use low-stakes self-assessment checklists that let the student record progress over time to visibly track mastery. Offer targeted praise focused on effort and strategy ("You used ratio reasoning correctly; that shows strong thinking") rather than only correctness. Include choice in practice activities (e.g., choose between a visual or algebraic representation) to build agency and preference expression. \newline 2) Concentrated errors on scaling word problems \newline Value: This pattern shows a specific gap in transferring proportional reasoning to contextual tasks; addressing it has high instructional value because it unlocks broader competence in real-world applications. \newline \textit{[\ldots\ truncated \ldots]} \newline Recommended supports: Lower the affective cost of reflection by providing sentence stems and brief, compelling reflection formats (e.g., one-sentence "I was surprised that..." or a quick two-choice poll about confidence). Integrate peer-model reflections (examples from classmates or teacher) to show the value of reflection. Offer the student a choice of reflection modalities (typed sentence, audio note, or drawing) so they can express preference and feel ownership. Set short, achievable reflection goals (e.g., submit one reflection per week) and celebrate completion to build valuation of the practice. \newline Overall prioritization and integration: I would first address low confidence because it affects willingness to attempt harder tasks; next, target the specific gap in scaling word problems with contextualized practice; concurrently, lower the barrier to reflection to cultivate the student's valuing of metacognition. For each recommendation, I would explicitly connect tasks to the student's interests and offer choices to strengthen their sense of ownership---framing practices as decisions the student makes rather than requirements---to align instructional moves with affective growth as well as cognitive skill-building. \\
  \midrule
  \textbf{Kimi-K2.6} & score 10/10 \\
  & The user wants me to act as a teaching assistant and analyze how to prioritize instructional strategies and motivational supports for a middle-school math student learning proportional relationships, specifically using the affective domain at the "Valuing" level (Krathwohl's taxonomy). \newline Key requirements: \newline 1. Identify three specific learner signals (from response patterns, question frequency, or demonstrated understanding) \newline 2. Explain the value assigned to each signal in shaping the learning path (connecting to Valuing level) \newline 3. Describe concrete instructional strategies and supports for each signal to foster engagement and ownership \newline The affective domain levels (Krathwohl): Receiving, Responding, Valuing, Organization, Characterization. At the Valuing level, the learner attaches worth and value to particular objects, phenomena, or behaviors. They demonstrate this by showing preference, seeking out activities, etc. So I need to frame the instructional strategies around helping the student *value* the learning process, the mathematical content, and their own agency. \newline \textit{[\ldots\ truncated \ldots]} \newline Wait, is Signal 1 "demonstrated understanding"? Yes, it's the demonstration of procedural understanding without conceptual. Signal 2 is question frequency. Signal 3 is response patterns (engagement/effort). Great. \newline Let me write the final response. \\
  \textbf{Innospark-1T} & score 8/10 \\
  & Let me think through this carefully. \newline The question/task: You are a teaching assistant supporting an instructor who will receive a personalized learner summary and suggested learning path based on analysis of an individual student's response patterns, question frequency, and demonstrated understanding in a middle-school math unit on proportional relationships. Using the affective (valuing) dom \newline My reasoning: Instructional Strategies Aligned with Affective Valuing Domain \newline Analysis Framework \newline The affective domain's Valuing level (Level 3 in Krathwohl's taxonomy) involves students attaching worth to particular objects, phenomena, or behaviors---moving beyond mere willingness to respond toward genuine commitment and preference. At this stage, students begin to find value in learning activities and develop internal motivation. Three Learner Signals and Prioritization \newline Signal 1: Quality of Explanations in Written Responses \newline \textit{[\ldots\ truncated \ldots]} \newline This prioritization aligns with affective valuing principles by: \newline 1. Progressing from external supports toward internal motivation \\
  \specialrule{1.2pt}{8pt}{8pt}
  \multicolumn{2}{@{}l}{\textbf{120730} \quad $\Delta = -20\%$} \\
  \textbf{Question} & As a research assistant supporting an educational experiment design, explain how you would use an automated statistical analysis pipeline (powered by AI) to (a) recommend group assignments and variable definitions, (b) preprocess and clean the collected data, and (c) generate and interpret the primary statistical results. In your explanation, describe each step you would expect the AI to perform, the rationale for those steps, and one example of an appropriate statistical test the pipeline might choose for a randomized two-group pretest-posttest design. \\
  \textbf{Reference} & To support an educational experiment design using an automated statistical analysis pipeline, I would organize the workflow into three main stages---design recommendations, data preprocessing, and analysis \& interpretation---and specify what I expect the AI to do at each stage and why. \newline (a) Recommend group assignments and variable definitions \newline - Steps expected from the AI: \newline 1. Clarify the research question and identify the primary outcome variable (e.g., student test score improvement) and key covariates (e.g., pretest score, grade level, prior achievement). \newline 2. Propose group assignment strategies that match the study goals and constraints, such as simple random assignment for internal validity or stratified randomization to balance important covariates (e.g., stratify by pretest score tertiles or school). \newline 3. Define operational variable measures: specify how the outcome is computed (posttest minus pretest or percent change), coding for categorical variables, and units for continuous variables. \newline \textit{[\ldots\ truncated \ldots]} \newline - The ANCOVA estimate indicates the treatment group scored on average 4.2 points higher on the posttest than the control group after adjusting for pretest scores (95\% CI: 1.1 to 7.3; p = 0.008). The standardized effect size (Cohen's d) is 0.35, a small-to-moderate educationally meaningful improvement. Model diagnostics show approximately normal residuals and no strong heteroscedasticity, and sensitivity analyses using robust regression produce similar effect estimates. Limitations include 8\% attrition balanced across groups and that the sample is from two schools, which may limit generalizability. \newline Overall, the AI-powered pipeline should provide reproducible, documented steps that assist the researcher from design through interpretation: recommend principled randomization and clear variable definitions, perform systematic cleaning and variable construction, select appropriate models (e.g., ANCOVA for pretest-posttest randomized designs), check assumptions, and produce clear numeric and plain-language summaries of the findings along with potential limitations. \\
  \midrule
  \textbf{Kimi-K2.6} & score 10/10 \\
  & The user wants an explanation of how an automated statistical analysis pipeline (powered by AI) would be used in three specific areas for an educational experiment design: \newline (a) recommend group assignments and variable definitions \newline (b) preprocess and clean the collected data \newline (c) generate and interpret the primary statistical results \newline For each step, I need to describe: \newline 1. What the AI would do \newline \textit{[\ldots\ truncated \ldots]} \newline - One example of appropriate statistical test: ANCOVA, yes \newline Looks good. I'll write it concisely. \\
  \textbf{Innospark-1T} & score 8/10 \\
  & Let me think through this carefully. \newline The question/task: As a research assistant supporting an educational experiment design, explain how you would use an automated statistical analysis pipeline (powered by AI) to (a) recommend group assignments and variable definitions, (b) preprocess and clean the collected data, and (c) generate and interpret the primary statistical results. In your expl \newline My reasoning: Automated Statistical Analysis Pipeline for Educational Experiments \newline Overview \newline As a research assistant, I would leverage an AI-powered statistical analysis pipeline to streamline the experimental design and analysis process. Below is a detailed explanation of how this pipeline would be applied across three critical phases of the research process. (a) Recommending Group Assignments and Variable Definitions \newline Steps the AI Would Perform: \newline \textit{[\ldots\ truncated \ldots]} \newline - Offering interpretive guidance grounded in statistical best practices \newline The automated approach doesn't replace researcher judgment---it augments it by handling technical complexities while flagging decisions requiring human oversight. \\
  \specialrule{1.2pt}{8pt}{8pt}
  \multicolumn{2}{@{}l}{\textbf{17305} \quad $\Delta = -20\%$} \\
  \textbf{Question} & Using Gagn's ABCD (Audience, Behavior, Condition, Degree) format and his principles for writing learning objectives, design one affective-domain learning objective at the Organizing level (prioritizing values and resolving conflicts) for a K--12 health education unit on promoting respectful peer communication about mental health. Then, explain how each ABCD component maps to Gagn's principles, justify why this objective targets the affective Organizing level, and describe one classroom assessment task and one instructional strategy aligned to that objective. \\
  \textbf{Reference} & Learning objective (ABCD): Given a small-group classroom discussion scenario and a rubric outlining respectful communication behaviors (Condition), a 10th grade student (Audience) will evaluate two conflicting peer attitudes about disclosing mental health struggles, prioritize the values involved (e.g., empathy, privacy, honesty), and produce a written action plan that reconciles those values by selecting and justifying at least three communication choices consistent with respectful peer support (Behavior), achieving a clear, logically ordered plan with at least 3 justified choices and rationale scored at 4 out of 5 on the rubric (Degree). \newline Mapping to Gagn's principles: \newline - Audience: The objective specifies the learner (10th grade student). Gagn emphasizes identifying who will perform the learning so instruction can be tailored to prior skills and developmental level. Specifying grade level ensures appropriate cognitive and social maturity for affective organizing tasks. \newline - Behavior: The stated behavior---evaluate conflicting attitudes, prioritize values, and produce a written action plan selecting and justifying communication choices---is observable and internalizes the affective process into a concrete product (a plan and justifications). Gagn requires behaviors to be stated as learning outcomes that can be observed or measured; here the written plan and justifications provide evidence of internal value reorganization. \newline - Condition: The condition names the instructional context and supports (small-group discussion scenario and rubric). Gagn notes that conditions describe what materials, prompts, or constraints will be provided; specifying the scenario and rubric clarifies the expected context for performance. \newline - Degree: The objective sets criteria---a logically ordered plan with at least three justified choices and a rubric score of 4/5---providing a measurable standard for acceptable performance, which aligns with Gagn's need for clear performance criteria. \newline \textit{[\ldots\ truncated \ldots]} \newline Instructional strategy aligned to the objective: \newline Strategy: Structured value reconciliation protocol. Steps: (1) Present the scenario and clarify the rubric; (2) Conduct a jigsaw discussion where students first list values independently, then meet in mixed groups to compare lists and debate tensions between values; (3) Guide students through a graphic organizer that asks them to rank values, note potential conflicts, and propose communication options that honor the top-ranked values; (4) Students draft their action plans and peer-review each other's plans using the rubric, focusing on whether the selected communication choices reflect the prioritized values. This strategy uses social interaction and scaffolded reflection to prompt internalization and reorganization of values, supports Gagn's sequence of events (gaining attention with a realistic scenario, informing learners of objectives and criteria, stimulating recall of prior attitudes, presenting guidance, eliciting performance via the plan, and providing feedback through peer review), and directly prepares learners for the assessment task. \\
  \midrule
  \textbf{Kimi-K2.6} & score 10/10 \\
  & [9-10] The objective is very specific, clear, and measurable. It explicitly describes the expected learner behaviors using appropriate action verbs (identify, prioritize, resolve, propose) and directly aligns with the affective organizing level and the associated assessment task. \\
  \textbf{Innospark-1T} & score 8/10 \\
  & Let me work through this systematically. \newline **Understanding Gagn's ABCD Format:** \newline - Audience: Who will be doing the learning \newline - Behavior: What observable action the learner will demonstrate \newline - Condition: Under what circumstances/context \newline - Degree: The standard/criteria for acceptable performance \newline \textit{[\ldots\ truncated \ldots]} \newline These must align with the Organizing level - students need to demonstrate they can prioritize values and resolve conflicts. \newline Let me draft the complete response. \\
  \specialrule{1.2pt}{8pt}{8pt}
  \multicolumn{2}{@{}l}{\textbf{109355} \quad $\Delta = -20\%$} \\
  \textbf{Question} & You are evaluating a high-school student's multi-paragraph essay about local watershed conservation. The student's essay contains these three sample sentences (excerpts). For each sentence: (1) identify the specific error or weakness (e.g., factual inaccuracy, unsupported claim, faulty reasoning, ambiguous wording, omission of evidence, inappropriate generalization); (2) explain why it undermines the argument or misleads the reader; and (3) propose a concrete, evidence-based revision or addition that corrects the problem and strengthens the essay. Evaluate all three excerpts and then, based on your analysis, provide a brief (2--3 paragraph) personalized explanatory paragraph the teacher could give the student that synthesizes the revisions and suggests next steps for improving the overall essay. \newline Excerpts: \newline A. "The local river has always been clean, so recent industrial development in the valley cannot be the cause of the fish die-offs reported this summer." \newline B. "Phosphate runoff is bad because it makes the water green, and the government should ban all fertilizers to protect aquatic life." \newline C. "A study I read said invasive plants reduce native biodiversity by up to 90\%, so planting native species around the river will instantly restore the original ecosystem." \newline Provide a detailed, evaluative response for each excerpt and the final personalized explanatory paragraph. \\
  \textbf{Reference} & Excerpt A --- Identification of error or weakness: The sentence contains a hasty generalization and an unsupported causal inference. It assumes a stable historical baseline ("has always been clean") without evidence and dismisses recent industrial development as a possible cause despite temporal correlation with fish die-offs. It also ignores alternative causes and lacks citation of water-quality data. \newline Why it undermines the argument: By asserting an absolute and unsupported historical claim, the sentence weakens credibility. The categorical dismissal of industrial development as a cause prevents the essay from evaluating plausible hypotheses and biases the reader against investigating relevant evidence. This reasoning could lead the reader to overlook correlated or contributing factors (e.g., chemical spills, thermal pollution, combined impacts with agricultural runoff, or delayed ecological responses). \newline Concrete, evidence-based revision: Revise to acknowledge uncertainty, cite data, and recommend investigation. Example: "Historical monitoring records indicate that the river's measured pollutant concentrations were relatively low between 2000 and 2015 (Smith et al., 2018); however, the recent timing of reported fish die-offs following new industrial activity in the valley suggests industrial sources should be investigated alongside other possible causes such as agricultural runoff and elevated water temperatures. Conducting targeted chemical analyses and comparing current water-quality metrics to baseline records would clarify likely causes." \newline Excerpt B --- Identification of error or weakness: This sentence contains an oversimplified causal claim, ambiguous terminology, and an unqualified policy prescription. It conflates visible algal growth ("makes the water green") with ecological harm and leaps to a sweeping policy solution ("ban all fertilizers") without assessing practicality, proportionality, or unintended consequences. \newline Why it undermines the argument: The lack of specificity (which phosphates, what concentrations, where) and absence of evidence (linking local fertilizer use to measured phosphate levels and ecological harm) makes the claim unconvincing. The blanket policy recommendation is extreme and may be counterproductive (e.g., harming agricultural productivity or prompting illegal use) and fails to consider targeted alternatives such as buffer strips, controlled application timing, or best-management practices. \newline Concrete, evidence-based revision: Offer a measured, evidence-linked claim and feasible policy options. Example: "Excess phosphate runoff from nearby agricultural fields can stimulate algal blooms that reduce dissolved oxygen and stress aquatic organisms; local water sampling in June found phosphate concentrations above the state guidance level of 0.1 mg/L (Environmental Agency, 2023). Rather than an outright ban, targeted strategies---such as establishing vegetated riparian buffers, promoting precision fertilizer application, and incentivizing reduced phosphate formulations---would lower runoff while balancing agricultural needs." \newline \textit{[\ldots\ truncated \ldots]} \newline Personalized explanatory paragraph for the student (2--3 short paragraphs): Your essay raises several important concerns about the river and shows promising engagement with sources, but some sentences make absolute claims or dramatic leaps from evidence to conclusion. To strengthen your argument, avoid blanket statements and provide context: cite baseline monitoring data before declaring that a river has "always" been clean, qualify study results by reporting the conditions under which they apply, and link causal claims to specific local measurements (e.g., phosphate concentrations, documented timing of industrial discharges). When proposing solutions, favor feasible, evidence-based strategies and explain how they would be implemented and evaluated rather than offering immediate or total bans. \newline For revisions, replace unsupported absolutes with nuanced, evidence-linked statements and recommend concrete next steps (targeted sampling, comparative data analysis, pilot restoration projects, and monitoring). This approach will make your essay more persuasive by showing you can weigh evidence, consider alternative explanations, and design realistic interventions---skills that align with environmental science reasoning and will strengthen both your claims and your recommendations. \\
  \midrule
  \textbf{Kimi-K2.6} & score 10/10 \\
  & [9-10] The response accurately identifies specific errors for all three excerpts, clearly explains how each undermines the argument with evidence, and provides concrete, evidence-based revisions. The personalized paragraph effectively synthesizes the feedback and offers tailored, actionable next steps (source audit, stakeholder paragraph, realistic timelines) perfectly suited to the learner's level. \\
  \textbf{Innospark-1T} & score 8/10 \\
  & Let me work through this problem step by step. \newline First, I need to understand what's being asked: You are evaluating a high-school student's multi-paragraph essay about local watershed conservation. The student's essay contains these three sample sentences (excerpts). For each sentence: (1) identify the specific error or weakness (e.g., factual inaccuracy, unsupported claim, faulty reasoning, ambiguous wording, omission of evidence, inappropriate generalization); (2) expl \newline Key values given: ,, ,, ,, 2, 3, 90, , \newline My approach: \newline 1. sentence presents a logical fallacy known as "false cause" or "post hoc ergo propter hoc" (after this, therefore because of this). The student assumes that because the river was historically clean, industrial development cannot be the cause of recent fish die-offs. This is faulty rea \newline 2. sentence contains multiple weaknesses: an oversimplified causal claim, an unsupported policy recommendation, and imprecise language. The student claims phosphate runoff "makes the water green" without explaining the actual mechanism (algal blooms), and jumps to an extreme policy conclusion ("ban all fertili \newline 3. sentence presents a "miracle cure" fallacy. The student cites a study (without proper attribution) claiming invasive plants reduce biodiversity "by up to 90\%," then makes an unsupported leap that planting native species will "instantly restore" the ecosystem. This ignores ecological complexity, time \newline This gives us the final answer: (after this, therefore because of this). The student assumes that because the river was historically clean, industrial \\
  \specialrule{1.2pt}{8pt}{8pt}
  \multicolumn{2}{@{}l}{\textbf{8678} \quad $\Delta = -20\%$} \\
  \textbf{Question} & Describe how you would carry out a needs analysis focused on clarifying a problem situation and defining the scope for a higher-education business administration project. In your description, list the key steps you would take to confirm the problem with stakeholders, gather and organize primary and secondary data, perform a basic cause-and-effect analysis to identify root causes, and explain how you would determine and delimit the scope of the solution. Also indicate what attitudes or stakeholder willingness (affective aspects) are important during the initial conversations and why they matter. \newline Give the final answer directly, without reasoning, and do not use the reasoning task such as think\textgreater{}. \\
  \textbf{Reference} & To carry out a needs analysis that clarifies the problem situation and defines the scope for a business-administration project, I would follow a sequence of clear steps combining stakeholder conversations, data collection, and causal analysis, while maintaining an openness to stakeholder perspectives. First, I would initiate stakeholder conversations to confirm the problem: schedule short, focused interviews with key stakeholders (e.g., department heads, instructors, students, administrators) and use structured questions to elicit their perceptions of the problem, its impacts, and desired outcomes. During these conversations I would adopt an attitude of active listening, neutrality, and willingness to hear differing views so participants feel heard and provide candid information. This affective openness helps surface underlying issues and builds trust for later data collection. \newline Second, I would gather and organize primary and secondary data. Primary data would include direct evidence collected specifically for the analysis: interview notes, focus-group summaries, surveys of affected users, classroom observations, and performance metrics (e.g., course completion rates, student satisfaction scores). Secondary data would include existing reports, institutional records, prior evaluation results, policy documents, and relevant literature. I would create a simple data inventory and categorize sources by type, date, and relevance, ensuring transparency about data limitations. \newline Third, I would synthesize the collected information and perform a basic cause-and-effect (Ishikawa/fishbone) analysis to identify possible root causes. I would begin by stating the observed problem clearly (e.g., declining course completion in a required course), then map contributing factors across categories such as People (instructor expertise, student preparedness), Process (curriculum alignment, assessment practices), Materials (instructional resources, technology), and Environment (class size, scheduling). For each branch I would note supporting evidence from the primary and secondary data. Where multiple plausible causes exist, I would prioritize them by frequency and impact using simple criteria (how often a factor appears across sources and how strongly it correlates with the problem). \newline Fourth, based on the cause analysis, I would determine and delimit the scope of the solution. I would define scope boundaries by specifying what will be addressed (e.g., redesign of course assessment and instructor support) and what is out of scope (e.g., institution-wide policy changes that require separate approval). Scope decisions would be guided by feasibility, stakeholder priorities uncovered in conversations, and available resources. I would document goals, target population, timeline, and constraints so all stakeholders share a common understanding of the project limits. \newline Finally, I would present a concise needs-analysis report summarizing the confirmed problem, supporting data, identified root causes, recommended focus areas, and the proposed scope with rationale. Throughout, I would emphasize maintaining an open, receptive stance with stakeholders---being willing to revise assumptions when new evidence appears---because this affective willingness to listen ensures the analysis reflects real needs and secures stakeholder buy-in for proposed solutions. \\
  \midrule
  \textbf{Kimi-K2.6} & score 10/10 \\
  & [9-10] The student fully and accurately recognizes the problem situation, clearly explaining the context. The output thoroughly covers problem confirmation, data gathering, cause-and-effect analysis, scope definition, and affective aspects. It is highly relevant and detailed enough to allow immediate progression to the next step of requirement analysis. \\
  \textbf{Innospark-1T} & score 8/10 \\
  & Let me think through this carefully. \newline The question/task: Describe how you would carry out a needs analysis focused on clarifying a problem situation and defining the scope for a higher-education business administration project. In your description, list the key steps you would take to confirm the problem with stakeholders, gather and organize primary and secondary data, perform a basic cause-and-effect analysis to ide \newline My reasoning: Needs Analysis for Higher-Education Business Administration Project \newline Overview \newline A systematic needs analysis establishes the foundation for project success by ensuring the right problem is addressed with appropriate boundaries. Below is a comprehensive approach to conducting this analysis. Key Steps to Confirm the Problem with Stakeholders \newline 1. Conduct Stakeholder Identification and Mapping \newline \textit{[\ldots\ truncated \ldots]} \newline 4. Stakeholders are engaged and committed to the process \newline The investment in thorough needs analysis pays dividends throughout the project lifecycle by reducing scope creep, preventing solution mismatch, and building the stakeholder buy-in essential for implementation success. \\
  \specialrule{1.2pt}{8pt}{8pt}
  \multicolumn{2}{@{}l}{\textbf{46528} \quad $\Delta = -20\%$} \\
  \textbf{Question} & You are a teaching assistant supporting kindergarten science teachers. Given a dataset of 120 learner questions collected from a unit on plants (questions include 'Why do leaves change color?', 'How do seeds grow?', 'My plant died---what went wrong?', repeated simple factual queries, and several misconceptions like 'Plants get food from soil directly' and 'Only sunlight makes plants grow fast'), analyze the question patterns and misconceptions. Break down the dataset to identify at least four distinct question categories (e.g., factual recall, procedural/help-seeking, conceptual misconception, curiosity/extension), describe the defining features and indicators of each category, estimate the relative frequency of each category with justification, and recommend three targeted feedback strategies or in-class interventions that the instructor can use to address the most common misconception and reduce repetitive low-value questions. In your analysis, explain how you would use these patterns to adapt upcoming lesson plans and formative assessment prompts for kindergarten learners. \\
  \textbf{Reference} & Analysis of the 120 learner questions: Category breakdown and indicators \newline 1) Conceptual misconception questions (approx. 40\% --- \textasciitilde{}48 questions) \newline - Defining features: Questions that reveal incorrect underlying ideas about plant biology rather than missing facts. Examples: 'Plants get food from soil directly,Only sunlight makes plants grow fast.' These questions often include incorrect causal statements or confuse processes (nutrition vs. water uptake vs. sunlight). They may be phrased confidently or as why/how but reflect an incorrect model. \newline - Indicators in dataset: Presence of assertions framed as questions, repeated wording suggesting the same false belief across multiple students (e.g., many variants about soil as food). \newline 2) Factual recall questions (approx. 25\% --- \textasciitilde{}30 questions) \newline - Defining features: Requests for single facts students lack (e.g., 'What do plants need to grow?', 'Where do seeds come from?'). Short, specific, and often asked early in unit or after demonstrations. \newline \textit{[\ldots\ truncated \ldots]} \newline Overall instructional insight \newline - The dominance of conceptual misconceptions indicates the need to prioritize conceptual-change instruction before or alongside procedural activities. Procedural scaffolds reduce repetitive low-level queries and restore teacher capacity for addressing deeper misunderstandings. Regular, simple formative tasks (drawings, card sorts, picture choices) will let the instructor track progress and adapt instruction effectively. \\
  \midrule
  \textbf{Kimi-K2.6} & score 10/10 \\
  & [9-10] The question classification is highly accurate and well-justified, matching the four required categories. The proposed interventions are specific, immediately actionable, and developmentally appropriate for kindergarteners (e.g., tactile sorting, visual self-service board). Lesson adaptations and formative assessments directly target misconceptions and reduce repetitive queries, demonstrating exceptional educational validity. \\
  \textbf{Innospark-1T} & score 8/10 \\
  & Let me think through this carefully. \newline The question/task: You are a teaching assistant supporting kindergarten science teachers. Given a dataset of 120 learner questions collected from a unit on plants (questions include 'Why do leaves change color?', 'How do seeds grow?', 'My plant died---what went wrong?', repeated simple factual queries, and several misconceptions like 'Plants get food from soil directl \newline My reasoning: Analysis of Kindergarten Plant Science Questions \newline Dataset Breakdown and Categorization \newline Based on the provided dataset of 120 learner questions, I have identified four distinct question categories: \newline 1. Curiosity/Extension Questions (20\% - \textasciitilde{}24 questions) \newline \textit{[\ldots\ truncated \ldots]} \newline | Curiosity/Extension | \textasciitilde{}24 | "Why/What if" questions | Wonder Walls | \newline This analysis provides a comprehensive framework for using question patterns to drive responsive, developmentally appropriate science instruction for kindergarten learners. \\
  \specialrule{1.2pt}{8pt}{8pt}
  \multicolumn{2}{@{}l}{\textbf{2966} \quad $\Delta = -20\%$} \\
  \textbf{Question} & You are an instructional designer conducting a task analysis for a kindergarten-level social studies/civics instructional goal. The goal is: "Children will demonstrate respect for classroom rules and each other's rights during group activities." Using Gagn's learning taxonomy, classify this instructional goal into the appropriate learning domain and subtypes (e.g., verbal information, intellectual skills, cognitive strategies, motor skills, attitudes). Then, identify the main instructional steps (main events of instruction and sub-tasks) required to achieve this goal for kindergarten learners. Finally, critically justify your classification and the chosen main steps: explain why each element of Gagn's taxonomy applies or does not apply, how the main steps align with children's developmental abilities, and how these steps will facilitate transfer and retention of respectful behavior in group contexts. \\
  \textbf{Reference} & Classification: This instructional goal primarily fits within Gagn's "attitudes" category (a subset of the affective side in some interpretations) and also engages "intellectual skills" and "verbal information." Specifically: - Attitudes: The core aim---demonstrating respect for rules and peers---is an attitude/valuing outcome because it concerns dispositions, choices, and socially guided behavior. - Intellectual skills: Elements such as recognizing rules, applying turn-taking sequences, and solving simple social problems involve intellectual skills (rule application and discrimination among social actions). - Verbal information: Children must learn and recall explicit classroom rules, names for emotions and rights (e.g., "my turn,safe hands"), and teacher-provided social scripts. Motor skills are peripheral (e.g., using quiet hands) but not primary; cognitive strategies (self-monitoring) may be introduced at a simple level for older kindergarteners. Main instructional steps (aligned to Gagn's events and task sub-tasks): 1) Gain attention \& motivate: Use a short, engaging social story puppet skit showing respectful vs. disrespectful group play. 2) Inform learners of the objective: Clearly state the goal in child-friendly language ("We will practice being kind and following rules so everyone can play together"). 3) Stimulate recall of prior learning: Ask children to share times they felt happy when someone shared or sad when someone was interrupted, priming existing social experiences. 4) Present content (verbal information + modeling): Present 3--4 explicit classroom rules and simple social scripts; teacher and puppets model correct behaviors and phrases (e.g., "Can I have a turn?Thank you for waiting"). 5) Provide guided practice with feedback (intellectual skills application): Conduct structured role-plays in small groups where children rehearse scripts and rule-following; teacher gives immediate, specific feedback focused on observable behaviors. 6) Elicit performance under varied conditions (transfer practice): Create short play stations that vary materials, peer partners, and minor distractions so children practice demonstrating respect across contexts. 7) Provide reinforcement and formative assessment: Use positive reinforcement (praise, labeled tokens) and brief teacher-noted checklists to acknowledge observed respectful acts. 8) Promote retention and transfer (cognitive strategies \& attitudes): Introduce simple self-monitoring prompts (e.g., a visual "hands/voice" cue cards) and classroom routines that embed the behaviors daily; involve families with a take-home story so children practice at home. 9) Enhance generalization: Periodically change group configurations and contexts (art, snack, outdoor play) and ask children to identify how rules apply in each setting. Justification and critique: - Why Attitudes: The desired outcome requires internalizing a disposition toward others---deciding to respect rules and rights---so attitude is the primary classification. Attitudes are best developed through repeated modeled behavior, social reinforcement, and opportunities to choose respectful actions; the instructional steps emphasize modeling, practice, and reinforcement accordingly. - Why Intellectual Skills and Verbal Information are included: Children must recognize and apply rules (intellectual skills) and remember specific phrases and rule names (verbal information). Without these cognitive components, attitude-focused strategies would lack concrete behaviors to enact. Modeling and role-play target intellectual skill acquisition, while clear rule presentation addresses verbal information needs. - Developmental alignment: Kindergarteners learn well through multimodal, concrete, and social experiences. Puppetry and role-play provide concrete, emotionally engaging exemplars. Short, structured practice rounds respect limited attention spans. Visual cues and routines scaffold working memory limits and support predictable practice. Self-monitoring is simplified to visual prompts appropriate for young children. - Facilitating transfer and retention: Varied practice contexts, daily routines, family involvement, and explicit cues promote both retention (through repetition and retrieval practice) and transfer (by varying conditions and encouraging reflection on applicability). Reinforcement shapes consistent behavior, while brief formative checks signal progress and maintain motivation. - Potential limitations and mitigations: Class-wide reinforcement risks overgeneralizing praise; using labeled, specific feedback mitigates ambiguity. Some children may need individualized supports (e.g., social stories with fewer steps or role-play modeling repeated more times); planning tiered practice addresses varying developmental levels. In summary, classifying the goal primarily as an attitude outcome, supplemented by intellectual skills and verbal information, is justified because respectful group behavior requires a disposition plus concrete rule knowledge and skills. The main steps---attention, explicit objectives, modeling, guided practice, varied application, reinforcement, and generalization---are chosen to align with kindergarten cognitive and social development and to maximize retention and transfer of respectful behaviors. \\
  \midrule
  \textbf{Kimi-K2.6} & score 10/10 \\
  & [9-10] The student correctly identifies Attitude as the primary Gagn domain with appropriate supporting domains, perfectly aligning with the learner's level. Main steps are logically ordered, mapped to Gagn's events, and use observable verbs. The tree structure is complete, and the justification thoroughly addresses taxonomy, development, and transfer. \\
  \textbf{Innospark-1T} & score 8/10 \\
  & Let me work through this problem step by step. \newline First, I need to understand what's being asked: You are an instructional designer conducting a task analysis for a kindergarten-level social studies/civics instructional goal. The goal is: "Children will demonstrate respect for classroom rules and each other's rights during group activities." Using Gagn's learning taxonomy, classify this instructional goal into the appropriate learning domain and subtypes (e.g., ve \newline Key values given: ,, ,, ,, ,, ,, , \newline My approach: \newline 1. of Instruction for Kindergarten \newline Step 1: Gain Attention \newline \textit{[\ldots\ truncated \ldots]} \newline 7. of respectful behavior \newline This instructional design creates a comprehensive framework that respects kindergarten developmental levels while systematically building the attitudes and skills necessary for respectful group participation. \\
  \specialrule{1.2pt}{8pt}{8pt}
  \multicolumn{2}{@{}l}{\textbf{111006} \quad $\Delta = -20\%$} \\
  \textbf{Question} & As a research assistant using an AI tool that automatically checks for privacy, research ethics, data inconsistencies, and generates transparency reports for AI usage, describe what ethical concerns you would be most alert to when preparing a dataset for analysis. In your answer, list at least four specific concerns, briefly explain why each is important to research ethics and participant welfare, and describe one concrete action you would take to address each concern. \\
  \textbf{Reference} & I would be particularly alert to the following ethical concerns when preparing a dataset for analysis using such an AI tool: \newline 1) Informed consent and scope of data use: It is essential to ensure that participants consented to the types of analyses and secondary uses planned. This protects participant autonomy and legal compliance. Concrete action: Review consent forms and metadata to confirm permitted uses; if the intended analysis falls outside consented purposes, seek re-consent or limit the dataset to records with appropriate consent. \newline 2) Identifiability and re-identification risk: Even de-identified data can sometimes be re-identified through combinations of quasi-identifiers. This threatens privacy and confidentiality. Concrete action: Apply and verify de-identification techniques (e.g., remove direct identifiers, generalize dates/locations, assess k-anonymity) and ask the AI tool to run a re-identification risk assessment, then further redact or aggregate fields as needed. \newline 3) Data quality and bias/inaccuracy: Erroneous, missing, or non-representative data can produce misleading results and unfair impacts on groups. Ensuring accuracy supports integrity and justice. Concrete action: Use the AI's automated checks to detect inconsistent values, outliers, and missingness patterns; document corrections and, where appropriate, flag or exclude low-quality records and note potential biases in the dataset. \newline 4) Minimization and necessity: Collecting or using more personal data than necessary increases risk without benefit. Ethical practice requires using the minimum data needed to achieve research aims. Concrete action: Review variables and remove fields not required for the research question; configure the AI to generate a transparency statement explaining which data elements were kept and why. \newline 5) Transparency and accountability in AI processing: Participants and stakeholders should be informed about AI involvement in data handling and decision-making to maintain trust. Concrete action: Use the tool's automated transparency-report feature to produce a plain-language summary of what AI processes were applied, parameter settings, and human oversight steps, and include that summary in study documentation. \newline 6) Secure storage and access controls: Improper storage or broad access increases risk of unauthorized disclosure. Concrete action: Ensure data are stored in encrypted environments with role-based access, log access events, and configure the AI tool to respect these access controls when processing data. \newline By remaining alert to these concerns and taking the specified concrete actions---verifying consent scope, reducing re-identification risk, correcting data quality issues, minimizing data collected, documenting AI processing transparently, and enforcing secure access---I would help uphold participant welfare, legal and ethical standards, and research integrity. \\
  \midrule
  \textbf{Kimi-K2.6} & score 10/10 \\
  & [9-10] Comprehensively identifies four relevant ethical issues with clear explanations tied to participant welfare and research integrity. Cites specific principles/regulations (k-anonymity, GDPR, IRB) and provides systematic, feasible, and concrete actions for each. \\
  \textbf{Innospark-1T} & score 8/10 \\
  & Let me think through this carefully. \newline The question/task: As a research assistant using an AI tool that automatically checks for privacy, research ethics, data inconsistencies, and generates transparency reports for AI usage, describe what ethical concerns you would be most alert to when preparing a dataset for analysis. In your answer, list at least four specific concerns, briefly explain why each is importa \newline My reasoning: Ethical Concerns in Dataset Preparation: A Research Assistant's Guide \newline Thinking Process \newline When preparing datasets for analysis, research assistants must navigate complex ethical considerations that protect participants and ensure research integrity. Let me identify four critical concerns and develop practical responses for each. Four Critical Ethical Concerns in Dataset Preparation \newline 1. Privacy and Re-identification Risk \newline \textit{[\ldots\ truncated \ldots]} \newline Conclusion \newline These four concerns---privacy protection, representational fairness, documentation transparency, and regulatory compliance---form the foundation of ethical dataset preparation. By implementing the concrete actions described, research assistants can uphold ethical standards while generating the transparency reports that accountability in AI-assisted research demands. \\
  \specialrule{1.2pt}{8pt}{8pt}
  \multicolumn{2}{@{}l}{\textbf{20496} \quad $\Delta = -10\%$} \\
  \textbf{Question} & Using Gagn's ABCD (Audience, Behavior, Condition, Degree) model and the results of a task analysis for a higher education political science module on comparative electoral systems, design a single constructed-response assessment item that aligns with one specific instructional objective from the module. Describe: (1) the instructional objective stated in ABCD format including the learning content; (2) the assessment prompt (the item) that students will answer; and (3) how the item requires students to apply knowledge from the task analysis (explain which component tasks or cognitive steps from the analysis the student must perform to complete the item). Your response should produce a ready-to-use assessment item and a clear explanation of its alignment with the instructional objective and task analysis steps. \\
  \textbf{Reference} & Instructional objective (ABCD + content): \newline Audience: Given graduate-level political science students enrolled in a course on comparative electoral systems, \newline Behavior: students will analyze and rank three different mixed-member electoral system designs according to their likely effects on party system fragmentation and district-level accountability, \newline Condition: using a provided country case brief (demographic profile, party cleavages, existing district magnitude, and sample ballot formulas) and without external resources, \newline Degree: and justify their ranking with at least three evidence-based reasons per system, citing how each reason links to features in the case brief. \newline Learning content: mixed-member electoral systems (variations in list compensation, district magnitude, and thresholds) and theoretical mechanisms linking system design to party fragmentation and accountability. \newline \textit{[\ldots\ truncated \ldots]} \newline Task analysis linkage (what cognitive steps were applied): \newline To complete this item I applied these task-analysis-derived steps: (1) identify relevant case-brief features (societal cleavages, district magnitudes, turnout trends); (2) map system design features (compensation rule, threshold, list openness, vote structure, seat shares) to theoretical mechanisms affecting fragmentation and accountability; (3) evaluate how each mechanism would operate in Republic X by combining features from steps 1 and 2; (4) weigh trade-offs and rank options based on comparative impact on fragmentation versus accountability; and (5) justify the ranking with at least three linked, evidence-based reasons per option. These steps show direct application of task analysis outcomes to produce an aligned assessment response. \\
  \midrule
  \textbf{Kimi-K2.6} & score 8/10 \\
  & [7-8] The student's answer mostly aligns with the performance objectives and is actionable/observable, with a well-structured ABCD objective, detailed prompt, and clear task analysis alignment. However, it lacks the explicit scoring criteria (rubrics and correct answers) required for the highest band (9-10). \\
  \textbf{Innospark-1T} & score 7/10 \\
  & Let me work through this problem step by step. \newline First, I need to understand what's being asked: Using Gagn's ABCD (Audience, Behavior, Condition, Degree) model and the results of a task analysis for a higher education political science module on comparative electoral systems, design a single constructed-response assessment item that aligns with one specific instructional objective from the module. Describe: (1) the instructional objective stated in ABCD format incl \newline Key values given: ,, ,, ,, 1, 2, 3 \newline Working through the calculation: \textless{}analysis\textgreater{} \newline The user wants a constructed-response assessment item for a higher education political science module on comparative electoral systems. The request requires using Gagn's ABCD model for the instructional objective and linking the assessment to a task analysis. Since the specific task analysis results aren't provided, I will generate a plausible task analysis for a specific learning objective within this domain to demonstrate the process. Instructional Objective (ABCD Format) \newline Audience (A): Undergraduate political science students in a comparative politics course. Behavior (B): Analyze and compare the potential outcomes of different electoral systems. Condition (C): Given a hypothetical country profile (population distribution, party ideologies, historical voting patterns) and a specific electoral reform proposal. Degree (D): With 90\% accuracy in identifying the correct theoretical outcomes and providing a rationale based on Duverger's Law and district magnitude principles. Task Analysis \newline \textit{[\ldots\ truncated \ldots]} \newline Step 1 (Identify Formula): The student must distinguish between FPTP (Plurality) and List PR (Proportional). Step 2 (Analyze District Magnitude): The student must recognize that moving from single-member districts (Magnitude 1) to a national list (High Magnitude) changes the translation of votes to seats. Step 3 (Calculate/Estimate): The student must perform the mathematical conversion or logical estimation of how votes translate to seats in both systems (e.g., "Party A got 40\% of the vote but won 60\% of seats under FPTP; under PR they would get 40\% of seats"). Step 4 (Evaluate Consequences): The student must discuss the shift from a majoritarian "manufactured majority" to a fragmented parliament requiring coalitions, citing the 5\% threshold as a filter for extremism or small parties. \newline This gives us the final answer: s. \\
  \specialrule{1.2pt}{8pt}{8pt}
  \multicolumn{2}{@{}l}{\textbf{32743} \quad $\Delta = -10\%$} \\
  \textbf{Question} & You are creating prompts to generate non-text multimodal learning materials (images, audio, or short videos) for a kindergarten-level science lesson about the life cycle of a butterfly. List and describe at least five specific prompt elements or instructions you would include for each modality (image, audio, and short video) to ensure the generated materials are age-appropriate, scientifically accurate, and easy for young children to understand. For each element, briefly explain why it is important for kindergarten learners. \\
  \textbf{Reference} & Image prompts: \newline 1. Simple, bold visuals of each life stage (egg, caterpillar, chrysalis, butterfly) with clear labels: Ensures children can visually recognize and name each stage. \newline 2. Bright, high-contrast colors and large shapes with minimal background detail: Keeps focus on the concept and makes images accessible for young viewers. \newline 3. Friendly, realistic-but-cute character style (soft edges, smiling faces on characters optional) to engage preschoolers without anthropomorphizing away the science: Balances engagement with accurate representation. \newline 4. Consistent character or color scheme across stages (same leaf pattern or same caterpillar markings) to show continuity: Helps children understand these stages belong to the same organism. \newline 5. Include a simple directional flow or arrows between stages and one-sentence caption per stage using simple vocabulary: Makes the sequence clear and supports early reading skills. \newline \textit{[\ldots\ truncated \ldots]} \newline 5. Include a 5--10 second pause after the full cycle with a visual prompt (e.g., a question card "Which stage comes next?") and on-screen labels for children to point to: Encourages recall and interactive viewing. \newline Why these elements matter: Each element is chosen to match kindergarten cognitive abilities---simple vocabulary, repetition, clear sequencing, multimodal reinforcement (visual + audio), and engagement techniques---so children can remember and name the life stages accurately while remaining interested and emotionally comfortable with the content. \\
  \midrule
  \textbf{Kimi-K2.6} & score 8/10 \\
  & [7-8] The student clearly presents learning objectives, target audience, and modality-specific instructions with strong pedagogical justifications. However, it lacks some detailed format specifications (e.g., exact resolution, file format, precise audio/video length) and explicit safety/accessibility guidelines needed for immediate reproducibility, placing it in the 7-8 band. \\
  \textbf{Innospark-1T} & score 7/10 \\
  & Let me think through this carefully. \newline The question/task: You are creating prompts to generate non-text multimodal learning materials (images, audio, or short videos) for a kindergarten-level science lesson about the life cycle of a butterfly. List and describe at least five specific prompt elements or instructions you would include for each modality (image, audio, and short video) to ensure the generated materials are age-appropri \newline My reasoning: Creating Effective Multimodal Prompts for Kindergarten Butterfly Life Cycle Lessons \newline Overview \newline When generating AI prompts for kindergarten educational materials, we must balance scientific accuracy with developmental appropriateness. Below are detailed prompt elements for each modality, with explanations of their pedagogical importance. Image Prompt Elements \newline | Element | Description | Why It Matters for Kindergarten | \newline \textit{[\ldots\ truncated \ldots]} \newline 1. Slow, Smooth Transitions | "Use gradual morphing animations between stages lasting 3-5 seconds each" | Abrupt changes confuse young viewers. Slow transitions help children mentally track the transformation process. | \newline | 2. Minimal On-Screen Text | "Limit text to 2-3 simple words per scene; rely primarily on visuals and narration" | Most kindergarteners are pre-readers or early readers. | \\
  \specialrule{1.2pt}{8pt}{8pt}
  \multicolumn{2}{@{}l}{\textbf{97499} \quad $\Delta = -10\%$} \\
  \textbf{Question} & You are an evaluator conducting a summative assessment of a middle-school civics project on community budgeting. Using the rubric provided (criteria: Understanding of civic concepts, Use of evidence, Collaboration, Presentation clarity, and Reflection), analyze the strengths and weaknesses revealed by a student's final report and portfolio. For each rubric criterion: (1) identify specific evidence from the student's work that demonstrates a strength or a weakness, (2) explain how that evidence connects to program-level outcomes for the civic education unit, and (3) recommend one targeted instructional or program change to address the weakness or build on the strength. Conclude with an overall judgment about the program's effectiveness based on patterns across the rubric criteria and suggest one priority action for program improvement. \\
  \textbf{Reference} & Strengths and weaknesses by rubric criterion, with evidence, program connections, and recommendations: \newline 1. Understanding of civic concepts \newline - Evidence (Strength): The student's report accurately defines key terms such as "public good,tax allocation," and "stakeholder engagement," and explains how a city budget reflects competing community priorities. The student links historical examples (e.g., municipal responses to infrastructure needs) to current budget choices. \newline - Connection to program outcomes: This demonstrates the program outcome of conceptual understanding---students can explain civic concepts and apply them to local contexts, indicating the unit effectively taught foundational ideas and their relevance. \newline - Recommendation: To deepen transfer, incorporate more scaffolds that require students to apply concepts to multiple, contrasting municipalities (e.g., large city vs. small town) so students practice adapting definitions to varied contexts. \newline 2. Use of evidence \newline \textit{[\ldots\ truncated \ldots]} \newline - Overall effectiveness: The program effectively fosters civic understanding and collaboration, achieving intended learning in those areas. However, its effectiveness is limited by insufficient emphasis on data literacy and presentation skills, which reduces the real-world impact of student work and weakens the credibility of their recommendations. \newline - Priority action for improvement: Implement a unit-long scaffold on evidence use and communication that includes: (a) a short data-literacy module (reading budgets, basic calculations, sourcing), (b) iterative checkpoints where students must submit and receive feedback on data appendices, and (c) dedicated rehearsal sessions for presentations with focus on visual clarity. This single, integrated change targets the primary weakness revealed across rubric criteria and is likely to raise the overall quality and persuasiveness of student outputs. \\
  \midrule
  \textbf{Kimi-K2.6} & score 10/10 \\
  & The user wants me to act as an evaluator conducting a summative assessment of a middle-school civics project on community budgeting. However, the user hasn't provided the actual student work or final report/portfolio. I need to create a realistic, illustrative example of how such an evaluation would look, using hypothetical but authentic evidence from a student's work. \newline I need to structure the response around 5 rubric criteria: \newline 1. Understanding of civic concepts \newline 2. Use of evidence \newline 3. Collaboration \newline 4. Presentation clarity \newline \textit{[\ldots\ truncated \ldots]} \newline One more check: The user asked for "specific evidence from the student's work." My invented evidence is specific enough: "\$50,000 expenditure,19\% vs 44\%,rewrote introduction overnight," etc. This should satisfy the requirement. \newline I think I'm ready to write. \\
  \textbf{Innospark-1T} & score 9/10 \\
  & Let me think through this carefully. \newline The question/task: You are an evaluator conducting a summative assessment of a middle-school civics project on community budgeting. Using the rubric provided (criteria: Understanding of civic concepts, Use of evidence, Collaboration, Presentation clarity, and Reflection), analyze the strengths and weaknesses revealed by a student's final report and portfolio. For each rubric criterio \newline My reasoning: \textless{}Thinking\textgreater{} \newline I need to analyze this middle-school civics project comprehensively using the provided rubric. Let me work through each criterion systematically, identifying specific evidence, connecting to program outcomes, and making targeted recommendations. \textless{}/Thinking\textgreater{} \newline Evaluative Analysis of Middle-School Civics Project on Community Budgeting \newline Criterion 1: Understanding of Civic Concepts \newline \textit{[\ldots\ truncated \ldots]} \newline 3. Reflection prompts specifically targeting systemic analysis ("How do economic structures influence budget decisions?") \newline This targeted intervention addresses the root cause---insufficient scaffolding for higher-order thinking---rather than treating symptoms in individual students. \\
  \specialrule{1.2pt}{8pt}{8pt}
  \multicolumn{2}{@{}l}{\textbf{95261} \quad $\Delta = -10\%$} \\
  \textbf{Question} & Using the rubric dimensions of Organization, Evidence and Support, Clarity of Argument, and Use of Sources, explain how you would evaluate an essay's overall performance. For each rubric dimension, describe what to look for in the student's writing and provide one concrete example of evidence from an essay that would justify a high score for that dimension. \\
  \textbf{Reference} & To evaluate an essay's overall performance using the four rubric dimensions---Organization, Evidence and Support, Clarity of Argument, and Use of Sources---I would assess each dimension separately, noting specific observable features in the writing and citing concrete examples that justify a high score. Below I describe what to look for and give an example of evidence indicating strong performance for each dimension. \newline 1) Organization \newline What to look for: A clear overall structure with an engaging introduction that states the main claim, logically ordered paragraphs each focused on a single idea, effective transitions between paragraphs, and a concise conclusion that synthesizes main points. Paragraphs should have topic sentences and follow a coherent progression from one idea to the next. \newline Example of evidence for a high score: The essay opens with a brief hook and a thesis statement: "Implementing project-based learning in middle school technology classes increases student engagement and problem-solving skills." Each body paragraph begins with a topic sentence (e.g., "First, project-based learning fosters engagement through real-world relevance."), presents a focused point, and ends with a transition such as "Consequently, students are more motivated to apply learned skills," which links to the next paragraph about skill development. The conclusion restates the thesis and synthesizes the three main benefits discussed, showing clear organizational control. \newline 2) Evidence and Support \newline What to look for: Claims are supported by relevant, specific evidence---data, examples, or observations---that directly connect to the writer's claims. The essay differentiates between fact and opinion and uses evidence to strengthen reasoning rather than relying on vague assertions. \newline \textit{[\ldots\ truncated \ldots]} \newline Example of evidence for a high score: The essay references a peer-reviewed article and an educational policy report: "Smith et al. (2018) found increased engagement in STEM electives; the Department of Education (2020) recommends project-based approaches for middle school curricula." The student paraphrases Smith's findings accurately and uses a short, contextualized quotation from the policy report to highlight a recommendation. Citations are consistently formatted in-text, and the writer comments on differences in sample populations between sources, demonstrating critical integration rather than mere listing. \newline Putting it together: After evaluating each dimension with the described indicators and examples, I would synthesize these observations to form an overall judgment. An essay that exhibits clear organization, provides specific and relevant evidence, presents a logical and coherent argument, and integrates credible sources thoughtfully shows mastery at the Understanding level and provides strong justification for a high holistic rating. \\
  \midrule
  \textbf{Kimi-K2.6} & score 10/10 \\
  & The user wants an explanation of how to evaluate an essay using four specific rubric dimensions: Organization, Evidence and Support, Clarity of Argument, and Use of Sources. For each dimension, I need to: \newline 1. Describe what to look for in the student's writing \newline 2. Provide one concrete example of evidence from an essay that would justify a high score \newline I should keep it concise and accurate as per the system prompt. \newline Let me structure this clearly with four sections, one for each dimension. \newline **Organization**: Look for logical flow, effective transitions, clear introduction and conclusion, coherent paragraph structure. Example: A thesis statement at the end of the intro that previews main points, topic sentences that mirror the thesis structure, and transitions like "While X is important, Y presents an equally compelling consideration" that guide the reader between paragraphs. \newline \textit{[\ldots\ truncated \ldots]} \newline Use of Sources: Look for credible, relevant sources that are smoothly integrated, properly cited, and synthesized to advance the writer's own argument. Example: "The essay introduces a Harvard Business Review study with a signal phrase---'As Christensen and Raynor (2023) argue'---embeds a brief quotation about disruptive innovation, cites it in APA format, and then extends the authors' framework to analyze a contemporary tech startup rather than merely summarizing the source." \newline These look good. I'll present them concisely with clear headings or bullet points. \\
  \textbf{Innospark-1T} & score 9/10 \\
  & Let me think through this carefully. \newline The question/task: Using the rubric dimensions of Organization, Evidence and Support, Clarity of Argument, and Use of Sources, explain how you would evaluate an essay's overall performance. For each rubric dimension, describe what to look for in the student's writing and provide one concrete example of evidence from an essay that would justi \newline My reasoning: Evaluating an Essay Using Four Key Rubric Dimensions \newline Organization \newline What to Look For: \newline - Logical flow of ideas from introduction to conclusion \newline \textit{[\ldots\ truncated \ldots]} \newline 5. Consider the essay's purpose and context when determining final assessment \newline This comprehensive approach ensures that evaluation is systematic, transparent, and provides meaningful feedback for student improvement. \\
  \specialrule{1.2pt}{8pt}{8pt}
  \multicolumn{2}{@{}l}{\textbf{53778} \quad $\Delta = -10\%$} \\
  \textbf{Question} & As a tutor supporting a learner's understanding of technology and home economics concepts, describe how you would explain the importance of being open to new ideas and willing to learn (receiving) when using new kitchen technologies (e.g., smart appliances, recipe apps). Include specific examples of behaviors a learner could show to demonstrate they are receptive, and explain why this attitude helps them learn safely and effectively. \\
  \textbf{Reference} & Being open to new ideas and willing to learn (the affective level of receiving) means paying attention, showing interest, and accepting guidance when encountering new kitchen technologies. To explain this to a learner, I would start by defining the concept: being receptive involves noticing new information, listening to instructions, and being ready to try different tools or procedures without immediate resistance. \newline I would give concrete examples tied to kitchen technologies. For instance, when introduced to a smart oven with preset cooking modes, a receptive learner listens to the demonstration, reads the quick-start guide, and tries a preset before changing settings. With a recipe app, they allow the app to suggest ingredient substitutions or step-by-step timers rather than dismissing it outright. If a companion smart scale recommends portion adjustments, the learner considers the suggestion and tests it on a small batch. \newline Specific behaviors that show receptiveness include: asking clarifying questions (e.g., "How does this sensor change cooking time?"), following safety prompts (e.g., acknowledging alerts and pausing to read them), taking notes on unfamiliar features, practicing under supervision, and giving honest feedback about what they observed. They also show willingness by experimenting carefully---using a simple recipe first, comparing results, and adjusting based on what they learned. \newline This attitude helps them learn safely and effectively because new technologies often have safety features or operating nuances that reduce risk only if users pay attention. Being receptive increases the chance they will notice warnings, understand correct usage, and adopt best practices (like cleaning sensors or updating software). It also shortens the learning curve: by trying recommended settings and following guidance, they build skill and confidence faster. Finally, openness fosters continuous improvement---when learners accept feedback and new techniques, they can incorporate better methods, save time, and achieve more consistent cooking results while minimizing hazards. \\
  \midrule
  \textbf{Kimi-K2.6} & score 10/10 \\
  & The user wants me to act as a tutor explaining the importance of being open to new ideas and willing to learn ("receiving") when using new kitchen technologies. I need to: \newline 1. Explain the importance of this attitude in the context of smart appliances and recipe apps \newline 2. Provide specific examples of behaviors demonstrating receptiveness \newline 3. Explain why this attitude helps with safe and effective learning \newline This is an educational question about technology and home economics. I should be concise, accurate, and supportive in tone. \newline Structure: \newline \textit{[\ldots\ truncated \ldots]} \newline - Conclusion: Encouraging summary \newline Keep it concise but comprehensive. Use accessible language appropriate for a tutor. \\
  \textbf{Innospark-1T} & score 9/10 \\
  & Let me think through this carefully. \newline The question/task: As a tutor supporting a learner's understanding of technology and home economics concepts, describe how you would explain the importance of being open to new ideas and willing to learn (receiving) when using new kitchen technologies (e.g., smart appliances, recipe apps). Include specific examples of behaviors a learner could show to demonstr \newline My reasoning: Explanation of Receptiveness in Learning Kitchen Technology \newline Understanding "Receiving" in the Learning Context \newline When I explain the concept of being open to new ideas and willing to learn (receiving) to a learner, I would start by defining it clearly: \newline "Receiving means approaching new kitchen technologies with a mindset of curiosity rather than fear or resistance. It's about acknowledging that while traditional cooking methods have value, modern tools can enhance your skills and efficiency when you allow yourself to learn how to use them properly." \newline \textit{[\ldots\ truncated \ldots]} \newline 3. Describe a potential safety issue that could arise from being unreceptive to learning about a smart appliance \newline This approach ensures the learner not only understands the concept but can apply it practically in their home economics practice. \\
  \specialrule{1.2pt}{8pt}{8pt}
  \bottomrule
\end{longtable}